\section{The higher-star bound}\label{sec:stars}

The successive incidence numbers give the integer
\[
d=\gcdop\{\lambda_j-\lambda_{j+1}:1\le j<t\}.
\]
It controls the exponent of the selected quotient.

\begin{theorem}\label{thm:hs}
With the notation above,
\[
d\,\Hk\ \le\ L\ \le\ \Hk .
\]
Consequently $\Delt(D)=\Hk/L$ is finite, is killed by $d$, is generated by at
most $v-1$ elements, and has order dividing $d^{\,v-1}$. There is a canonical
exact sequence
\[
0\to\Delt(D)\to\coker(A)\to\ZZ/k\to0,
\]
and, for the retained-apex matrix $A^{+}=\left[\begin{smallmatrix}\mathbf
1^{T}\\ A\end{smallmatrix}\right]$, a canonical equality
$\Delt(D)=\mathrm{Tor}\,\coker(A^{+})$; the splitting
$\coker(A^{+})\cong\ZZ\oplus\Delt(D)$ holds noncanonically.
\end{theorem}

\begin{proof}
Column weight $k$ gives $L\le\Hk$.

For the lower bound, let $u_R$ be the sum of
all block columns containing a fixed $j$-set $R$; coordinate counting gives
$u_R=\lambda_{j+1}\mathbf 1_P+(\lambda_j-\lambda_{j+1})\mathbf 1_R$. Hence, for
distinct points $x,y$ and an $(j-1)$-set $S$ disjoint from them,
\[
u_{S\cup\{x\}}-u_{S\cup\{y\}}=(\lambda_j-\lambda_{j+1})(e_x-e_y)\in L ,
\]
and B\'ezout over $1\le j<t$ yields $d(e_x-e_y)\in L$ for all $x,y$. Fix a block
$B$ and a point $o\in B$; since $d\,\mathbf 1_B\in L$,
\[
d\,\mathbf 1_B-\sum_{x\in B\setminus\{o\}}d(e_x-e_o)=dk\,e_o\in L .
\]
The vectors $\{e_x-e_o:x\ne o\}$ together with $k\,e_o$ generate $\Hk$, and both
$d(e_x-e_o)$ and $dk\,e_o$ lie in $L$, so $d\Hk\le L$. Thus $\Delt(D)$ is finite
and killed by $d$.

In $\Hk/L$ the block relation
$\mathbf 1_B=k\,e_o+\sum_{x\in B\setminus\{o\}}(e_x-e_o)$ vanishes, expressing the
class of $k\,e_o$ through the $v-1$ point-difference classes; hence $\Delt(D)$ is
a quotient of $(\ZZ/d)^{\,v-1}$, so it has at most $v-1$ generators and
$|\Delt(D)|$ divides $d^{\,v-1}$.

The coordinate-sum map
$\coker(A)\to\ZZ/k$, $[z]\mapsto\sum z\bmod k$, has kernel $\Hk/L=\Delt(D)$ and is
onto because $[e_x]\mapsto1$ generates $\ZZ/k$.

Finally the primitive map
$\ZZ\oplus\ZZ^P\to\ZZ$, $(a,z)\mapsto ka-\sum z$, has kernel identified
with $\Hk$ by $z\mapsto(\sum z/k,z)$. The columns of $A^+$ belong to that
kernel. Quotienting by them gives
$0\to\Delt(D)\to\coker(A^+)\to\ZZ\to0$, which splits noncanonically
and identifies $\Delt(D)$ canonically with the torsion subgroup.
\end{proof}

The star mechanism (the identity producing the coefficients
$\lambda_r-\lambda_{r+1}$) is a specialization of Wilson's higher-incidence
identity~\cite{Wilson1982}; see~\cite[Section 2.2, equation (2.1)]{KiermaierWassermann}
with $e=1$ and $f=j$. The universal group presented by $A$, and the $t=2$
point-difference identity, are in~\cite[Section 2]{CFP2017}. The gcd over
star levels and the retained-row passage above are elementary deductions.

\begin{corollary}\label{cor:d1}
If $d=1$ then $L=\Hk$ and $\Delt(D)=0$, i.e.\ $\coker(A)\cong\ZZ/k$ and the full
incidence Smith form is $1^{\,v-1}k$.
\end{corollary}
